\part{Complex and transcendental non-Archimedean theory}
\label{part:complex}

This part begins on a compact K\"ahler manifold.  Test curves and geodesic
rays lead to inverse-system constructions of transcendental
non-Archimedean metrics, their operations and traces, and their associated
Okounkov and Duistermaat--Heckman data.  In the projective setting,
multiplier ideals provide a valuative realization and an order-preserving
comparison with the trivially valued BJ theory.

The second half constructs the Mesquita--Piccione semivaluation space
\(X^\beth\) for a compact complex space.  Its dual-complex, model, and
finite-energy pluripotential theory is developed under the compact K\"ahler
hypotheses used by the cited sources; the more general ideal-spectrum
construction is not made to supply those K\"ahler conclusions by itself.
The envelope, orthogonality, and Calabi--Yau theory of
Mesquita--Piccione--Witt Nystr\"om and the applications to the cscK problem
are then treated on their stated domains.  The later
pseudoeffective comparison with the Darvas--Xia--Zhang construction belongs
to \Cref{part:comparisons}, because it requires a cross-theory domain audit
rather than a change of notation inside either construction.

\paragraph{Prerequisites and support.}
\label{front:part5-prerequisites-and-support}
Complex psh functions, non-pluripolar products, energy, and geodesic rays are
recalled in \Cref{chap:complex-prerequisites}; a fuller compact reference is
given in \Cref{app:complex-reference}.  The valuative chapters use
\Cref{chap:valuations,chap:skeleta,chap:topological-tools}, and the projective
comparison presupposes the homogeneous FS and filtration language of
\Cref{chap:bj-fs-pl-psh,chap:bj-homogeneity-bdivisors}.  The analytic
modification and dual-complex construction is developed in
\Cref{chap:mp-dual-complexes}; Bingener's theorem is identified there as a
load-bearing input rather than hidden inside an algebraic analogy.  Its exact
use is the proper finite-presentation dimension comparison over the excellent
singleton Stein compactum \(\{p\}\): it turns the terminal algebraic center
provided by the local Zariski lemma into a codimension-one analytic center.
That step precedes invariant divisorial comparison, divisorial density, and
the inverse limit of dual complexes.

\paragraph{The DXZ cluster.}
\label{front:part5-the-dxz-cluster}
Geodesic rays and test curves are related by Legendre duality in
\Cref{chap:rays-test-curves}.  Inverse systems over K\"ahler perturbations
define transcendental metrics in \Cref{chap:transcendental-metrics}; their
operations and trace maps are treated in \Cref{chap:operations-trace}, and
convex bodies and Duistermaat--Heckman measures in
\Cref{chap:transcendental-okounkov-dh}.  Valuative analytification and the
projective comparison with BJ are given in
\Cref{chap:valuative-analytification,chap:dxz-bj-comparison}.  The
class-free tropical spectrum of \Cref{chap:tropical-spectra-complex}
provides the transition to the second half.

\paragraph{The MP cluster.}
\label{front:part5-the-mp-cluster}
The complex-analytic algebraization needed for the dual-complex inverse limit
is established in \Cref{chap:mp-dual-complexes}.  Model and completed MP psh
functions, mixed measures, and energy are developed in \Cref{chap:mp-psh};
the variational Calabi--Yau theory is treated in \Cref{chap:mpwn-cy}, and the
stability applications in \Cref{chap:mpwn-k-stability}.  These chapters keep
K\"ahler classes, big pseudoeffective classes, and zero-volume inverse systems
distinct, and they label every point at which compactness, envelope
continuity, or algebraization is imported.

\paragraph{Reading routes.}
\label{front:part5-reading-routes}
The \emph{short route} to transcendental metrics is
\Cref{chap:complex-prerequisites,chap:rays-test-curves,chap:transcendental-metrics,chap:transcendental-okounkov-dh}.
The \emph{standard route} reads the part in order and uses
\Cref{chap:dxz-bj-comparison} and \Cref{chap:tropical-spectra-complex} as the
two transition chapters.  The \emph{reference route} adds
\Cref{app:complex-reference,app:convexity,app:ideals,app:intersection-bdivisors}
and follows every construction through \Cref{chap:mpwn-k-stability}; it is
the appropriate preparation for
\Cref{chap:pseff-mp-dxz,chap:ytd-bridge}.  A reader concerned only with the
projective BJ comparison may stop after \Cref{chap:dxz-bj-comparison}; a
reader concerned with general compact complex spaces may begin with
\Cref{chap:mp-dual-complexes} after reviewing
\Cref{chap:complex-prerequisites}, but in the non-K\"ahler case only the
constructions explicitly declared there remain available.

\paragraph{Main outcomes and interfaces.}
\label{front:part5-main-outcomes-and-interfaces}
The part supplies two complementary routes from complex data to
non-Archimedean objects.  The DXZ route packages rays and test curves into a
K\"ahler-perturbation inverse system and, on the projective common domain,
connects that system to BJ filtrations and FS functions.  The MP route builds
the semivaluation space and its model-theoretic pluripotential calculus from
analytic modifications and dual complexes.  Their common ideal-valuative
content, and the limits of extending it to pseudoeffective classes, are
analyzed in \Cref{chap:pseff-mp-dxz}.  These constructions are analogous to
BFJ model limits but are not identified with a Berkovich space over a
nontrivially valued field; the exact arrows are recorded in
\Cref{chap:master-dictionary}.

On the smooth projective positive-volume locus,
\Cref{thm:ok-mvob-metrics} proves the mixed partial-Okounkov-body formula for
big line bundles endowed with singular psh metrics of positive Cao
(equivalently, ideal-model) volume; \Cref{thm:ok-mvob-line-bundles} is its
minimal-singularity movable-
product specialization.  The proof was supplied by the author in an
unpublished manuscript, repaired at the convex-sandwich, model-envelope,
analytic-resolution, and finite-level degeneration interfaces, and verified
first as a repaired proof object and then at both dated theorem
boundaries by the production protocol recorded in \Cref{app:provenance}.
The metric formula and its line-bundle specialization have separate final
\texttt{correct}, zero-error, zero-gap receipts for those archived bytes;
the current chapter records its later source and proof review separately.
This verification does not
turn the author-supplied work in progress into published literature or create
a priority claim.  The supremum over valuations is not asserted to be
attained, and no zero-volume, transcendental-class, or singular-variety
extension is inferred.

\paragraph{Normalization and domain warnings.}
\label{front:part5-normalization-and-domain-warnings}
Complex potentials use squared Hermitian norms, whereas non-Archimedean
metric weights use the norm itself; comparison therefore carries the
factor-of-two convention fixed in \Cref{app:normalizations}.  Ideal values
must also be distinguished modulo pointwise equality from their
multiplier-ideal, or \(\mathcal I\)-, equivalence class.  Energy,
Monge--Amp\`ere mass, and probability normalization require the positivity
and finite-energy hypotheses stated in the relevant chapter.  In particular,
big non-nef classes and zero-volume classes are not treated as formal limits
of the K\"ahler probability theory.  The factor-\(n\) correction in MPWN25
has an independent \texttt{gpt-5.6-sol}/\texttt{xhigh} receipt, recorded in
\Cref{chap:mpwn-k-stability,app:provenance}.  The dual-complex boundary
repair in MPWN25 has likewise passed independent
\texttt{gpt-5.6-sol}/\texttt{xhigh} verification, as recorded in
\Cref{chap:mpwn-cy,app:provenance}.  The discrepancy shift in MP24 remains
visibly marked as a pending correction and is not treated as independently
verified.
